% The picture and the data structure are the same object seen twice: every wall
% implied on the left is one internal node on the right, and the parameters
% labelling that node are exactly the t, j and i the algorithm drew.
\begin{figure}[htbp]
  \centering
  \begin{subfigure}[b]{0.36\textwidth}\centering
    \begin{tikzpicture}[scale=0.78]
      \fill[cellgrey] (0,0) rectangle (4,-4);
      \draw[black!60] (0,0) rectangle (4,-4);
      \draw[black!20] (1,0) -- (1,-4); \draw[black!20] (3,0) -- (3,-4);
      \draw[black!20] (0,-1) -- (4,-1); \draw[black!20] (0,-2) -- (4,-2);
      \draw[black!20] (0,-3) -- (4,-3);
      % division A: the root, a vertical wall at j=1, door at row i=1
      \draw[wall segment] (2,0) -- (2,-1);
      \draw[wall segment] (2,-2) -- (2,-4);
      \draw[door] (2,-1) -- (2,-2);
      \node[wallred, font=\scriptsize\bfseries] at (2.35,0.28) {A};
      % division B: inside the left half, horizontal wall at j=1, door at column 0
      \draw[wall segment] (1,-2) -- (2,-2);
      \draw[door] (0,-2) -- (1,-2);
      \node[wallred, font=\scriptsize\bfseries] at (-0.32,-2) {B};
      % division C: inside the right half, horizontal wall at j=2, door at column 3
      \draw[wall segment] (2,-3) -- (3,-3);
      \draw[door] (3,-3) -- (4,-3);
      \node[wallred, font=\scriptsize\bfseries] at (4.32,-3) {C};
    \end{tikzpicture}
    \caption{The first three divisions of a $4\times4$ area.}
    \label{fig:region-tree-cuts}
  \end{subfigure}\hfill
  \begin{subfigure}[b]{0.60\textwidth}\centering
    \begin{tikzpicture}[scale=1.0, level distance=13mm,
      level 1/.style={sibling distance=42mm},
      level 2/.style={sibling distance=21mm}]
      \node[tree node] {A: $t{=}1,\,j{=}1,\,i{=}1$}
        child {node[tree node] {B: $t{=}0,\,j{=}1,\,i{=}0$}
          child {node[tree leaf] {$\region{0,0}{1,1}$}}
          child {node[tree leaf] {$\region{0,2}{1,3}$}}
        }
        child {node[tree node] {C: $t{=}0,\,j{=}2,\,i{=}3$}
          child {node[tree leaf] {$\region{2,0}{3,2}$}}
          child {node[tree leaf] {$\region{2,3}{3,3}$}}
        };
    \end{tikzpicture}
    \caption{The top of the region tree it builds. The four boxes at the bottom
    are regions still to be cut, not cells.}
    \label{fig:region-tree-tree}
  \end{subfigure}
  \caption{Every division is one internal node. The node records the
  orientation $t$, the wall position $j$ and the door position $i$ it drew, and
  its two children are the two regions it created; a leaf is a corridor. A tree
  with $L$ leaves has $L-1$ internal nodes, so of the $WH-1$ doors of
  Proposition~\ref{prop:spanning-tree}, $L-1$ are division doors and the rest
  are laid down inside the corridors.}
  \label{fig:region-tree}
\end{figure}
